%%%%%%%%%%%%%%%%%%%%%%%%%%%%%%%%%%%%%%%%%%%%%%%%%%%%%%%%%%%%%%
%% Rozdział 7: Platforma edukacyjna
%%%%%%%%%%%%%%%%%%%%%%%%%%%%%%%%%%%%%%%%%%%%%%%%%%%%%%%%%%%%%%

\section{Platforma edukacyjna}
\label{chap:educational-platform}

Oprócz platformy do anotacji opisanej
w rozdziale~\ref{chap:annotation-platform}, która wspiera gromadzenie danych
 do trenowania modeli, projekt dostarcza platformę edukacyjną udostępniającą
wytrenowane modele rozpoznawania i tłumaczenia
(rozdziały~\ref{chap:isolated-gloss-recognition}
 i~\ref{chap:sentence-level-translation}) użytkownikom końcowym za pośrednictwem
przystępnego interfejsu.

\todo{DO UZUPEŁNIENIA: Tematyka tego rozdziału nie została jeszcze uwzględniona
w dotychczasowych raportach z prac badawczych i programistycznych projektu.
Rozdział należy opracować po zaimplementowaniu platformy edukacyjnej. Poniższe
podrozdziały przedstawiają jego planowaną strukturę.}

\subsection{Architektura i stos technologiczny}
\label{sec:edu-platform-architecture}

\todo{DO UZUPEŁNIENIA: Należy opisać architekturę systemu i stos technologiczny
platformy edukacyjnej, zachowując taki sam poziom szczegółowości jak
w podrozdziale~\ref{sec:annotation-system-architecture} dotyczącym platformy
 do anotacji (frameworki frontendu i backendu, przechowywanie danych,
infrastruktura wdrożeniowa).}

\subsection{Integracja modeli}
\label{sec:edu-platform-model-integration}

\todo{DO UZUPEŁNIENIA: Należy opisać sposób udostępniania platformie
wytrenowanego modelu rozpoznawania izolowanych glos
(rozdział~\ref{chap:isolated-gloss-recognition}) lub całościowego potoku
przetwarzania (rozdział~\ref{chap:sentence-level-translation}), bądź obu tych
rozwiązań -- infrastrukturę udostępniania modeli i wykonywania wnioskowania,
kwestie związane z opóźnieniami oraz sposób przekazywania predykcji do frontendu.}

\subsection{Przegląd interfejsu użytkownika}
\label{sec:edu-platform-ui}

\todo{DO UZUPEŁNIENIA: Należy opisać interfejs użytkownika i przebieg interakcji
(np.\ rozpoznawanie na żywo z wykorzystaniem kamery internetowej, strukturę lekcji
 i quizów, informację zwrotną dla osoby uczącej się), a po zaimplementowaniu
zamieścić zrzuty ekranu głównych widoków. Uwaga: repozytorium projektu zawiera już
wersję demonstracyjną rozpoznawania na żywo z kamery internetowej, łączącą enkoder
Transformer z wyszukiwaniem za pomocą FAISS (commit ``live recognition
Transformer + FAISS''), co stanowi naturalny punkt wyjścia dla tego podrozdziału.}